\subsection{Host Level Classification}
\label{sec5.3: host results}
\noindent Host level prediction extends our conformal prediction framework from binary attributes to a multiclass task across 
$N = 54$ candidate bacterial hosts ($\mathcal{H} = \{h_1, \ldots, h_54\}$). This setting introduces statistical and dimensional 
challenges: heavy label imbalance, small per host calibration samples, and inference through a hierarchical two-stage cascade. 
Rather than attempting to fit a $2\,036$ dimensional global envelope across the joint feature space, we constructed class 
conditional envelopes which were calibrated independently for each candidate host $h \in \mathcal{H}$ on its dedicated $K_h$ 
dimensional score sub vectors $\mathbf{s}_h$.



\subsubsection{Experimental Protocol}
\label{sec5.3.1: host formulation}
\noindent We apply the split conformal framework of Section~\ref{sec3: methodology} to the reconstructed score vectors (\ref{eq: 
host nc transformation}) of the $N_{\text{tot}} = 15{,}638$ annotated phages across $54$ candidate hosts ($K_h \in [26, 40]$ per 
host; $K_{\text{tot}} = 2{,}036$). 
We again use a three way partition of the genome cluster data:
\begin{itemize}
    \item $D_{\text{train}}$ (Split $2$--$4$, $N_{\text{train}} = 9{,}101$): For each candidate host with sufficient training 
          data ($n_h \ge 3$), the infector phages are partitioned $50/50$ into shape-discovery ($S_1^{(h)}$, $\lfloor 
          n_h/2 \rfloor$) and size-scaling ($S_2^{(h)}$, $\lceil n_h/2 \rceil$) subsets to fit $E_h$ and compute $\hat{t}_h$.
    \item $D_{\text{val}}$ (Split $1$, $N_{\text{val}} = 2{,}593$): used for envelope hyperparameter tuning.
    \item $D_{\text{test}}$ (Split $0$, $N_{\text{test}} = 3{,}944$): held out and evaluated once after fixing all 
          parameters.
\end{itemize}
\noindent Unlike Gram type classification (Section~\ref{sec5.2: gram results}), host level scores use per model split 
assignments (Section~\ref{sec4.3: contamination fix}) to ensure the out of split validity across all five splits rather than 
restricting calibration to split $0$ alone. Missing entries ($\text{NA}$) are handled by each envelope geometry without data 
imputation

\vspace{0.15cm}
\noindent At test time, the Gram restriction $\hat{\mathcal{H}}(\mathbf{s})$ narrows the host space to $16$ for Gram positive 
and $38$ for Gram negative predictions. The cascaded prediction set $C_{\text{host}}(\mathbf{s})$ is evaluated using the per 
host inclusion rule of Equation~\eqref{eq: host-prediction-set}. The cascade coverage guarantee of Proposition~\ref{prop:cascade} 
applies directly, bounding theoretical miscoverage by $\le 2\alpha = 0.20$. If a test 
phage falls outside the calibrated envelopes of all candidate hosts ($C_{\text{host}}(\mathbf{s}) = \emptyset$), the empty set is 
resolved by selecting the candidate minimizing normalized nonconformity:
\begin{equation}
\label{eq: empty-set-resolution}
h_{\text{forced}}(\mathbf{s}) = \operatorname*{arg\,min}_{h \in \hat{\mathcal{H}}(\mathbf{s})} \frac{\tau_h(\mathbf{s}_h)}
{\hat{t}_h},
\end{equation}
which normalizes distances across different host geometries without defaulting to the uninformative full label set.


% % \begin{figure}[t]
% % \centering
% % \includegraphics[width=0.8\linewidth]{host_envelope.png}
% % \caption{Illustration of per host nonconformity envelopes.}
% % \label{fig: host-envelope}
% % \end{figure}

% \noindent Because $S_2^{(h)}$ only contains true infectors of host $h$, each envelope is Mondrian calibrated conditionally 
% \ref{sec2.3.2: split conformal}on $H = h$:
% \begin{equation}
% \mathbb{P}\big(\tau_h(\mathbf{s}_h) > \hat{t}_h \,\big|\, H = h\big) \le \alpha_h,
% \end{equation}
% where exact validity $\alpha_h = \alpha$ holds whenever $|S_2^{(h)}| \ge \frac{1-\alpha}{\alpha} = 9$.

% \vspace{0.15cm}
% \noindent Three structural properties justify using the per host envelopes rather than a single global $2\,036$ dimensional 
% envelope:
% \begin{enumerate}[label=(\roman*)]
%     \item \textbf{Coordinate Heterogeneity:} Different host genera are characterized by distinct subsets of functional dimensions 
%            ($K_h \in [26, 40]$), making a unified global coordinate system impossible
%     \item \textbf{Class Conditional Validity:} Fitting envelopes per host prevents dominant host classes from skewing coverage 
%            for rare classes.
%     \item \textbf{Variance Control:} Operating in $\mathbb{R}^{K_h}$ ($K_h \ll 2\,036$) reduces sample sparsity and stabilizes 
%            quantile estimation.
% \end{enumerate}

%%%%%%%%%%%%%%%%%%%%%%%%%%%%%%%%%%%%%%%%%%%%%%%%%%%%%%%%%%%%%%%%%%%%%%%%%%%%%%%%%%%%%%%%%%%%%%%%%%%%%%%%%

\subsubsection{Hyperparameter Selection on Validation Data}
\label{sec5.3.2: host tuning}

\noindent We optimized the hyperparameters on $D_{\text{val}}$ (Split 1) by minimizing the average set size subject to empirical 
coverage constraints at target level $1 - \alpha = 0.90$::
\begin{equation}
\label{eq: constrained-opt}
\theta^\star \in \arg\min_{\theta} \overline{|C|}_{\text{val}}(\theta) \quad \text{subject to} \quad \hat{c}_{\text{val}}(\theta) 
\ge 1 - \alpha.
\end{equation}
\noindent with ties or near constraint cases resolved by minimizing the coverage discrepancy $|\hat{c}_{\text{val}}(\theta) - 
(1 - \alpha)|$.

\vspace{0.15cm}
\noindent Validation under the unrestricted candidate space selected:
\begin{itemize}
    \item \textbf{Strip Binning ($N_B$):} Swept over $N_B \in \{3, 5, 8, 10, 15, 20\}$ with windowed monotonic smoothing ($w=2$). 
           $N_B = 8$ was selected, achieving $90.2\%$ validation coverage with an average prediction set size of $13.43$ 
           candidate hosts.
    \item \textbf{Radial Resolution ($M, \kappa$):} Swept across $M \in \{100, 250, 500\}$ directions and concentration parameters 
           $\kappa \in \{4.0, 8.0, 16.0\}$. $M = 250$ with $\kappa = 4.0$ provided the optimal trade-off closest to nominal 
           validity ($89.0\%$ validation coverage, average set size $12.75$).
\end{itemize} 


%%%%%%%%%%%%%%%%%%%%%%%%%%%%%%%%%%%%%%%%%%%%%%%%%%%%%%%%%%%%%%%%%%%%%%%%%%%%%%%%%%%%%%%%%%%%%%%%%%%%%%%%%

\subsubsection{Empirical Test Results (Split 0)}
\label{sec5.3.3: host test results}

\noindent Table~\ref{tab: host_results} summarizes performance metrics evaluated on the test split ($N_{\text{test}} = 3{,}944$ 
phages) across all 54 hosts.

\begin{table}[htbp]
\centering
\resizebox{\columnwidth}{!}{%
\begin{tabular}{lccccc}
\toprule
\textbf{Method} & \textbf{Coverage} & \textbf{Avg.\ Set Size} & \textbf{Singleton Rate} & \textbf{Singleton Acc.} & \textbf{Forced 
Rate} \\
\midrule
Radial Envelope ($M=250, \kappa=4$) & 0.933 & 15.74 & 0.066 & 0.762 & 0.016 \\
Strip Envelope ($N_B=8, w=2$)      & 0.904 & 15.55 & 0.001 & --    & 0.000 \\
Collapsed 1D                        & 0.927 & 7.11  & 0.166 & 0.907 & 0.006 \\
\bottomrule
\end{tabular}}
\caption{Performance metrics on the test split}
\label{tab: host_results}
\end{table}

\begin{figure}[htbp]
    \centering
    \includegraphics[width=1\textwidth]{Plots/Host coverage and set size.png}
    \caption{Host classification performance on Split 0 test phages ($1-\alpha = 0.90$). Left panel: overall empirical coverage 
            compared against the target (red dashed line). Right panel: average prediction set size across candidate hosts.}
    \label{fig: host coverage set size}
\end{figure}

\noindent As demonstrated in Table~\ref{tab: host_results} and Figure~\ref{fig: host coverage set size}, all three geometries 
achieve empirical coverage close to the nominal $90.0\%$ target. In contrast to the binary Gram setting where max threshold 
calibration yielded $99.8\%$ coverage, independent per host calibration eliminates artificial over coverage. 

\vspace{0.15cm}
\noindent Prediction efficiency diverges sharply:
\begin{enumerate}[label=(\roman*)]
    \item \textbf{Efficiency of Collapsed 1D:} Collapsed 1D achieves an average prediction set size of $7.11$, which is less than 
           half that of Radial ($15.74$) or Strip ($15.55$). In moderate to high dimensions ($K_h \ge 26$), 1D mean pooling avoids 
           the sample sparsity penalties that arise due to angular cones and multi dimensional bins. The complete per host set size 
           distributions and boxplots across all 54 hosts are in Appendix \ref{app: plots} 
           Figure~\ref{fig: set_size_distribution_histogram} and \ref{fig: set_size_distributions_boxplot}
    \item \textbf{Singleton Accuracy:} Collapsed 1D isolates a single host for $16.6\%$ of test phages and achieves $90.7\%$ 
           accuracy on these singletons. In contrast, Radial produces singletons for only $6.6\%$ of test phages (with only $76.2\%$ 
           accuracy), while Strip almost never gives a singleton output ($0.1\%$).
    \item \textbf{Forced Decision Rate:} The forced fallback rule \eqref{eq: empty-set-resolution} was only used for $0.6\%$ of 
           test phages for Collapsed 1D and $1.6\%$ for Radial. This confirms that host specific envelopes reliably capture 
           empirical infector score distributions.
\end{enumerate}


\noindent Figure~\ref{fig: per_host_coverage} in the Appendix \ref{app: plots} illustrates the class conditional validity across all 
hosts. The coverage remains well calibrated across major host types (e.g., \emph{Escherichia}, \emph{Pseudomonas}, 
\emph{Staphylococcus}, and \emph{Mycobacterium}), but we can also find low coverage to the host classes with sparse data.


%%%%%%%%%%%%%



% \subsubsection{Geometric Boundary Visualization}
% \label{sec5.3.4: geometric comparison}

% \noindent To illustrate the structural mechanisms driving these differences, Figure~\ref{fig: envelope geometry staph} displays the calibrated boundary projections on a two-dimensional subspace (lysin vs.\ annealing functional scores) for \emph{Staphylococcus}.

% \begin{figure}[htbp]
% \centering
% \includegraphics[width=\linewidth]{Plots/staph_lysin_annealing_comparison.png}
% \caption{Per-host envelope boundary comparison for \emph{Staphylococcus} across the three constructions, projected onto the (lysin, annealing) functional-score plane. Left: Collapsed 1D linear threshold. Centre: Radial smooth arc. Right: Strip piecewise-constant staircase.}
% \label{fig: envelope geometry staph}
% \end{figure}

% \noindent The geometric profiles illustrate clear qualitative distinctions:
% \begin{itemize}
%     \item \textbf{Collapsed 1D (Linear Cut):} Forms an anti-diagonal linear boundary governed by the mean nonconformity threshold. While unable to distinguish whether predictive signal originates from lysin or annealing, it avoids directional sparsity and yields compact sets in aggregate.
%     \item \textbf{Radial (Smooth Angular Arc):} Generates smooth arcs that curve inward for phages with balanced functional scores. However, because training scores concentrate along coordinate axes, the envelope expands along both axes, admitting more false positives in high dimensions.
%     \item \textbf{Strip (Piecewise Staircase):} Traces empirical density clusters via axis-aligned bins, achieving the lowest false-negative count at the expense of a conservative staircase boundary in data-sparse regions.
% \end{itemize}

% \subsubsection{Effect of Hierarchical Gram-Type Filtering}
% \label{sec5.3.5: gram filtering effect}

% \noindent Table~\ref{tab: gram filtering effect} compares Collapsed 1D performance with and without Gram-type pre-filtering.

% \begin{table}[htbp]
% \centering
% \caption{Impact of hierarchical Gram-type filtering on Collapsed 1D host prediction.}
% \label{tab: gram filtering effect}
% \begin{tabular}{lcc}
% \toprule
% \textbf{Configuration} & \textbf{Coverage} & \textbf{Avg.\ Set Size} \\
% \midrule
% Unrestricted (Full $\mathcal{H}$, $N=54$) & 0.925 & 7.12 \\
% Gram-Type Filtered ($\mathcal{H}_{\text{cand}}$) & 0.927 & 7.11 \\
% \bottomrule
% \end{tabular}
% \end{table}

% \noindent Gram-type pre-filtering produces a modest set-size reduction ($7.12 \rightarrow 7.11$). This demonstrates that per-host functional envelopes are already specific; non-target hosts of the opposing Gram phenotype are rejected natively by the host-specific nonconformity thresholds without relying heavily on the upstream taxonomic filter.

% \subsubsection{Robustness Ablations and Methodological Limitations}
% \label{sec5.3.6: host ablations}

% \noindent We conducted three targeted ablation experiments to evaluate specific operational boundaries:
% \begin{itemize}
%     \item \textbf{Low-Data Host Sensitivity:} Restricting evaluation to the 17 host genera with $\ge 100$ calibration phages increases Strip coverage from $0.897$ to $0.969$. This confirms that Strip's slight aggregate under-coverage is driven primarily by low-sample genera where bin marginals cannot be estimated robustly.
%     \item \textbf{Dimensionality Reduction:} Restricting host subspaces to their top-$3$, $5$, or $8$ most informative functional dimensions degraded empirical coverage to $86\text{--}88\%$. This indicates that secondary functional annotations provide discriminative signal rather than noise.
%     \item \textbf{Two-Dimensional Functional Subspaces:} Evaluating phages using only lysin and tail-fiber dimensions (available for 51 of 54 hosts) yielded coverages of $0.857$, $0.910$, and $0.928$ for Collapsed 1D, Radial, and Strip, respectively. Strip outperforms other geometries in this low-dimensional regime, where its 2D grid avoids the curse of dimensionality.
% \end{itemize}

% \paragraph{Documented Negative Result: Impostor-Bounded Calibration.}
% An alternative calibration strategy aimed at bounding false-positive rates by incorporating non-infector scores ("impostor calibration") was implemented. This formulation led to near-total prediction abstention (empty sets in $>99.9\%$ of test cases), resulting in $0\%$ empirical coverage across all geometries. Because this reflects an over-constrained optimization rather than a valid conformal set predictor, we document it as an open limitation to be addressed in future work (Section~\ref{sec6.3: future work}).


%%%%%%%%%%%%%%%%%%%%%%%%%%%%%%%%%%%%%%%%%%%%%%%%%%%%%%%%%%%%%%%%%%%%%%%%%%%



% \subsubsection{Geometric Envelope Comparison}
% \label{sec5.3.4: geometric comparison}

% \noindent To visualize the structural mechanics driving these performance differences, Figure~\ref{fig: envelope geometry staph} displays the calibrated boundary projections on a two-dimensional subspace (lysin vs.\ annealing functional scores) for \emph{Staphylococcus}.

% \begin{figure}[htbp]
% \centering
% \includegraphics[width=\linewidth]{Plots/staph_lysin_annealing_comparison.png}
% \caption{Per-host envelope boundary comparison for \emph{Staphylococcus} across the three constructions, projected onto the (lysin, annealing) functional-score plane. Left: Collapsed 1D linear threshold. Centre: Radial smooth arc. Right: Strip piecewise-constant staircase.}
% \label{fig: envelope geometry staph}
% \end{figure}

% \noindent The geometric profiles illustrate clear qualitative distinctions:
% \begin{itemize}
%     \item \textbf{Collapsed 1D (Linear Cut):} Collapsed 1D constructs an anti-diagonal linear boundary governed strictly by the mean nonconformity threshold. While blind to whether predictive evidence originates from lysin or annealing features, it produces compact boundaries in aggregate.
%     \item \textbf{Radial (Smooth Angular Arc):} Radial envelopes generate smooth arcs that curve inward for phages exhibiting balanced intermediate scores. However, because training scores concentrate along coordinate axes, the envelope expands along both axes, admitting more false positives in high dimensions.
%     \item \textbf{Strip (Piecewise Staircase):} Strip bounds closely trace empirical density clusters via axis-aligned bins, achieving the lowest false-negative count at the expense of a conservative staircase boundary in data-sparse regions.
% \end{itemize}

% \subsubsection{Effect of Hierarchical Gram-Type Filtering}
% \label{sec5.3.5: gram filtering effect}

% \noindent To assess the empirical benefit of the hierarchical Gram-type cascade, Table~\ref{tab: gram filtering effect} compares Collapsed 1D performance with and without Gram-type pre-filtering.

% \begin{table}[htbp]
% \centering
% \caption{Impact of hierarchical Gram-type filtering on Collapsed 1D host prediction.}
% \label{tab: gram filtering effect}
% \begin{tabular}{lcc}
% \toprule
% \textbf{Configuration} & \textbf{Coverage} & \textbf{Avg.\ Set Size} \\
% \midrule
% Unrestricted (Full $\mathcal{H}$, $N=54$) & 0.925 & 7.12 \\
% Gram-Type Filtered ($\mathcal{H}_{\text{cand}}$) & 0.927 & 7.11 \\
% \bottomrule
% \end{tabular}
% \end{table}

% \noindent Gram-type pre-filtering produces a minor improvement in set efficiency ($7.12 \rightarrow 7.11$). This demonstrates that the per-host functional envelopes are already highly specific; non-target hosts of the opposing Gram phenotype are rejected natively by the host-specific nonconformity thresholds without relying heavily on the upstream taxonomic filter.

% \subsubsection{Robustness Ablations and Methodological Limitations}
% \label{sec5.3.6: host ablations}

% \noindent We conducted three targeted ablation experiments to evaluate specific operational boundaries:
% \begin{itemize}
%     \item \textbf{Low-Data Host Sensitivity:} Restricting evaluation to the 17 host genera with $\ge 100$ calibration phages increases Strip coverage from $0.897$ to $0.969$. This confirms that Strip's slight aggregate under-coverage is driven primarily by low-sample genera where bin marginals cannot be estimated robustly.
%     \item \textbf{Dimensionality Reduction:} Restricting host subspaces to their top-$3$, $5$, or $8$ most informative functional dimensions degraded empirical coverage to $86\text{--}88\%$. This indicates that secondary functional annotations provide discriminative signal rather than noise.
%     \item \textbf{Two-Dimensional Functional Subspaces:} Evaluating phages using only lysin and tail-fiber dimensions (available for 51 of 54 hosts) yielded coverages of $0.857$, $0.910$, and $0.928$ for Collapsed 1D, Radial, and Strip, respectively. Strip outperforms other geometries in this low-dimensional regime, where its 2D grid avoids the curse of dimensionality.
% \end{itemize}

% \paragraph{Documented Negative Result: Impostor-Bounded Calibration.}
% An alternative calibration strategy aimed at bounding false-positive rates by incorporating non-infector scores ("impostor calibration") was implemented. This formulation led to near-total prediction abstention (empty sets in $>99.9\%$ of test cases), resulting in $0\%$ empirical coverage across all geometries. Because this reflects an over-constrained optimization rather than a valid conformal set predictor, we document it as an open limitation to be addressed in future work (Section~\ref{sec6.3: future work}).

%%%%%%%%%%%%%%%%%%%%%%%%%%%%%%%%%%%%%%%%%%%%%%%%%%%%%%%%%%%%%%%%%%%%%%%



% \subsubsection{Geometric Comparison Across Envelope Methods}
% \label{sec5.3.3: geometric comparison}

% \noindent Sections~\ref{sec5.3.2: host coverage} reports aggregate numbers; Figure~\ref{fig: envelope geometry 
% staph} makes the underlying geometric difference between the three constructions directly visible, using a 
% two-dimensional projection of \emph{Staphylococcus}'s own score space (lysin vs.\ annealing) with true/false 
% positive and negative test points overlaid.

% \begin{figure}[htbp]
% \centering
% % \includegraphics[width=\linewidth]{Plots/staph_lysin_annealing_comparison.png}
% \caption{Per-host envelope boundary comparison for \emph{Staphylococcus} across the three constructions, projected 
% onto the (lysin, annealing) functional-score plane. Collapsed 1D's boundary (left) is the straight anti-diagonal 
% line implied by its scalar mean-threshold (Section~\ref{sec3.4: envelope construction}); Radial's boundary (centre) 
% is the smooth arc produced by the von Mises–Fisher-weighted radius of Equation~\eqref{eq: radial boundary}; Strip's 
% boundary (right) is the piecewise-constant staircase produced by its binned conditional limits.}
% \label{fig: envelope geometry staph}
% \end{figure}

% \noindent The three panels make concrete what Section~\ref{sec3.4: envelope construction} states abstractly: 
% Collapsed 1D's envelope is a straight line, blind to whether a phage's evidence for \emph{Staphylococcus} comes from 
% lysin or annealing so long as their average clears the threshold — visible in its $21$ false positives 
% (\texttt{own\_cov}=$0.94$) sitting exactly on the wrong side of a linear cut that a curved boundary would exclude. 
% Radial's arc curves inward for phages strong in one dimension but weak in the other, at the cost of its own failure 
% mode: $2{,}568$ false positives, far more than Collapsed's $21$, because a boundary built from calibration data 
% concentrated near the axes extends generously along both axes individually. Strip's staircase sits between the two, 
% tracking the true cluster's shape more tightly than the line but not as smoothly as the arc, with the fewest false 
% negatives of the three ($5$, vs.\ Collapsed's $15$ and Radial's $10$) at the cost of the most false positives 
% ($2{,}540$).

% \vspace{0.15cm}
% \noindent This single host and dimension pair is illustrative, not exhaustive. \textbf{[TO BE COMPLETED ONCE YOU 
% GENERATE THE ADDITIONAL PANELS — see instructions below.]} We repeat this comparison for [$N$ additional] 
% functional-dimension pairs across [list hosts], confirming the same qualitative pattern: Collapsed 1D 
% under-adapts to the true cluster shape, Radial over-extends along calibration-dense directions, and Strip tracks the 
% cluster boundary most closely of the three at the cost of coarser, bin-driven steps.

% \subsubsection{Effect of Gram-Type Filtering}
% \label{sec5.3.4: gram filtering effect}

% \noindent To assess whether the gram-type cascade of Section~\ref{sec3.1: problem formalisation} materially improves 
% host-level prediction, we compare Collapsed 1D's performance with and without the gram-type restriction 
% $\hat{\mathcal{H}}(\mathbf{s})$:

% \begin{table}[htbp]
% \centering
% \caption{Effect of gram-type filtering, Collapsed 1D method.}
% \label{tab: gram filtering effect}
% \begin{tabular}{lcc}
% \toprule
% \textbf{Configuration} & \textbf{Coverage} & \textbf{Avg.\ Set Size} \\
% \midrule
% No gram-type filter (full $\mathcal{H}$) & 0.925 & 7.12 \\
% Gram-type filtered ($\hat{\mathcal{H}}(\mathbf{s})$) & 0.927 & 7.11 \\
% \bottomrule
% \end{tabular}
% \end{table}

% \noindent The difference is negligible for Collapsed 1D — filtering candidate hosts by predicted gram-type barely 
% moves either coverage or set size. This is a real, if modest, empirical result worth stating plainly rather than 
% assuming the cascade obviously helps efficiency: for this method, most of the discriminating power evidently comes 
% from the per-host envelopes themselves, not from gram-type pre-filtering. We note this comparison was only run for 
% Collapsed 1D; whether Radial and Strip show the same near-equivalence remains open, since a full comparison would 
% require the same baseline evaluation for both.

% \subsubsection{Robustness Ablations}
% \label{sec5.3.5: host ablations}

% \noindent We report three further ablations that probe specific failure modes of the main result above, rather than 
% attempting to improve on it directly.

% \vspace{0.15cm}
% \noindent \textbf{Low-data hosts.} Restricting evaluation to the $17$ hosts with $\ge 100$ training phages (of $54$ 
% total) raises Strip's coverage to $0.969$, well above its full-$\mathcal{H}$ result of $0.897$ — consistent with 
% Strip's per-bin quantile calibration (Section~\ref{sec3.4: envelope construction}) being the most sample-hungry of 
% the three constructions, and suggesting its below-target aggregate coverage in Table~\ref{tab: host results} is 
% driven disproportionately by hosts with too few calibration phages for Strip's bin structure to calibrate reliably.

% \vspace{0.15cm}
% \noindent \textbf{Feature selection.} Restricting each host to its top-$3$/$5$/$8$ most informative functional 
% dimensions, rather than the full $K_h$, \emph{reduces} coverage ($0.969 \to 0.86$–$0.88$) despite the lower 
% dimensionality. This is the opposite of what a naive curse-of-dimensionality argument would predict, and suggests 
% the additional dimensions dropped were carrying genuine discriminating signal rather than noise.

% \vspace{0.15cm}
% \noindent \textbf{Restricted functional categories.} Using only the lysin and tail-appendage columns (available for 
% $51$ of $54$ hosts) yields coverage $0.857$/$0.910$/$0.928$ for Collapsed/Radial/Strip respectively — below the 
% full-feature result for Collapsed and Radial, but \emph{above} it for Strip, which benefits from the sharply reduced 
% bin-count requirement of a two-dimensional problem.

% \subsubsection{A Negative Result: Specificity-Controlled Calibration}
% \label{sec5.3.6: specificity controlled}

% \noindent We attempted a specificity-controlled ("impostor-bounded") calibration variant, suggested by our 
% supervisor, intended to bound the false-positive rate against non-target hosts directly rather than only through 
% $\alpha$. As implemented, this experiment produced near-total abstention (empty prediction sets in 
% $99.9$–$100\%$ of cases) and $0\%$ empirical coverage across all three envelope constructions — not a valid 
% conformal outcome under Theorem~\ref{thm: marginal coverage}, and evidence of an implementation defect rather than a 
% genuine finding. We report this as a documented limitation: the specificity-controlled direction remains 
% unresolved and is left for future work (Section~\ref{sec6.3: future work}), rather than a completed result.